% ===================================================================
\chapter{The Baseline: Kiem's Thesis}
\label{ch:kiem}

Most results in this thesis are set beside a result from Lukas Kiem's
master's thesis \cite{kiem2025}, carried out at Infineon. This chapter sets out
what he did, on what data, and what he concluded. Section and table numbers
quoted here (``his Table~3.5'') refer to his thesis.

\section{His problem and his data}
\label{sec:kc_problem}

Kiem addressed the transport problem of Section~\ref{sec:motivation} with a
use case (his Section~3.3): one Infineon CTRX8191F MMIC with four receive
channels and 12-bit samples at \SI{50}{MS/s} produces
\begin{equation}
\label{eq:kc_rate}
50\,\mathrm{MS/s} \times 12\,\mathrm{bit} \times 4 = \SI{2.4}{Gbit/s},
\end{equation}
so one Gigabit Ethernet link needs a compression ratio of 2.4, or 2.7 with
protocol headers. The sensor cycles at \SI{20}{\hertz}; a frame of 1024 ramps
takes about \SI{21}{\milli\second} to acquire, and he allowed compression
\SI{15}{\milli\second} of the remaining \SI{29}{\milli\second}.

Two features of his setting differ from this thesis's. His radar uses
\textbf{one transmitter}, so consecutive ramps are consecutive observations by
the same antenna pair and the question on which this thesis's prediction design
turns (Section~\ref{sec:predaxis}) never arises. And it delivers \textbf{real}
samples, whose spectrum is symmetric, so he keeps half of each range FFT
because the other half is redundant; ColoRadar's samples are complex.

His evaluation chain (his Figure~3.1), which this thesis adopts almost
unchanged, shifts the ADC samples into a 32-bit word, applies a Hanning window
and a range FFT divided by the number of samples, compresses and decompresses,
and then runs a Doppler FFT, non-coherent integration (NCI) over the receive
channels and a peak detector, comparing the detections with those of an FP16
reference path. Every compression result he publishes, in simulation and in
hardware, was measured on his own recordings: one MMIC, $512 \times 512$
samples, four receive channels, 14-bit data. They are not public. He also
shows a synthetic five-target scene in white noise (his Table~3.2), but only
as an illustration of the chain.

\section{His design space and result}
\label{sec:kc_dse}

He compared FP16 (the reference), 16-bit fixed point (FX16), two Infineon
common-exponent formats (CMPRA and CMPRB, fixed ratios 2 and 4), Texas
Instruments' exponential-Golomb engine at four settings (EGE10 to EGE4, ratios
1.6 to 4), and Doppler Redundancy Hybrid Encoding (DRHE)
\cite{li2014,zivkovic2015}: prediction along the slow-time axis of the
range-FFT output with two IIR filters in polar form, then Huffman coding of the
differences. He excluded one method, the LPC-Huffman of Meucci and Mancuso
\cite{meucci2022}, in these words:

\begin{quote}\small
``LPC-Huffman compression is not taken because it has the same approach as
DRHE, but instead of compressing range FFT data, it operates on raw ADC data;
the results for DRHE are more promising.'' (his Section~3.4.1)
\end{quote}

\noindent
He assumed that method's domain from its int16 input, did not run it, and
compared published figures from different data. Table~\ref{tab:kc_t35} is his
result. DRHE reached a ratio of 3.25 with exactly FX16's detection metrics,
since it is lossless with respect to the 16-bit codes; the lossy schemes either
stayed below the requirement or lost detections. One miss in about 170
(\SI{0.59}{\percent}) shows that his recording held about 170 reference
detections, at an SNR of about \SI{22}{\decibel}.

\begin{table}[htbp]
\centering
\caption[Kiem's design-space results]{Kiem's design-space exploration on his
recordings (his Table~3.5). FP\% is normalised by test detections, so it is a
false discovery rate (Section~\ref{sec:detectionmetrics}).}
\label{tab:kc_t35}
\small
\begin{tabular}{@{}lrrrrr@{}}
\toprule
Algorithm & CR & FN\% & FP\% & NF (dBFS) & SNR (dB)\\
\midrule
FP16  & 1    & 0     & 0     & $-70.714$ & 22.407\\
FX16  & 1    & 0.59  & 1.18  & $-70.691$ & 22.330\\
CMPRA & 2    & 0.59  & 0.59  & $-70.687$ & 22.418\\
CMPRB & 4    & 7.69  & 4.73  & $-70.253$ & 22.111\\
EGE10 & 1.6  & 1.18  & 1.18  & $-70.689$ & 22.403\\
EGE8  & 2    & 2.37  & 2.37  & $-70.640$ & 22.289\\
EGE6  & 2.67 & 7.69  & 7.10  & $-70.367$ & 22.197\\
EGE4  & 4    & 40.24 & 26.63 & $-74.689$ & 28.245\\
\textbf{DRHE} & \textbf{3.25} & 0.59 & 1.18 & $-70.691$ & 22.330\\
\bottomrule
\end{tabular}
\end{table}

\section{His final scheme}
\label{sec:kc_adapt}

Two simplifications made DRHE suitable for hardware. He removed the run-length
half of the original coder, which changed the ratio only from 3.25 to 3.24 and
lets every value be coded independently in one clock. And he fixed the Huffman
dictionary, built once from the $S_4$ statistics of one of his recordings (his
Appendix~A); it is the \emph{original dictionary} of this thesis
(Section~\ref{sec:dict}). For each range bin $n$ of ramp $m$ and each channel,
his final scheme (his Section~3.5.3) predicts
\begin{align}
\hat{a}[m,n] &= \alpha\hat{a}[m-1,n] + (1-\alpha)\,a[m-1,n],
\label{eq:kc_mag}\\
\hat{\theta}[m,n] &= \beta\hat{\theta}[m-1,n] + (2-\beta)\,\theta[m-1,n] - \theta[m-2,n],
\label{eq:kc_phase}
\end{align}
with $\alpha = 0.6$ and $\beta = 0.4$, wraps the phase once into $[-\pi,\pi]$,
converts the prediction back to Cartesian with rounding, and codes the 16-bit
difference as its size category $S_4$ (Table~\ref{tab:catranges}) and the
APPEND bits. The state starts at zero and is updated from the input samples.
These are Equations~\eqref{eq:magpred} and~\eqref{eq:phasepred} of
Chapter~\ref{ch:methods} at a lag of one ramp.

\section{His hardware}
\label{sec:kc_hw}

He implemented the scheme with Vitis HLS and Vivado 2022.2 on an AMD Zynq
UltraScale+ \texttt{XCZU3CG-SFVC784-1-e}, in a system that streams
pre-windowed ADC data from an SD card through a Xilinx FFT core into the
compressor. The compressor reads one 128-bit word per cycle (one complex bin
for each of four channels), uses fixed point throughout, holds its state in
block-RAM arrays of 1024 words per channel, packs each channel's codes into a
58-bit word, and runs at $\mathrm{II} = 1$ at \SI{100}{\mega\hertz}.
Table~\ref{tab:kc_hwres} collects the figures he reports. His throughput is
theoretical, $128 \times \SI{100}{\mega\hertz}$ expressed in binary units; his
hardware ratio is below his simulated one because the hardware FFT is fixed
point.

\begin{table}[htbp]
\centering
\caption[Kiem's hardware results]{Kiem's published hardware figures (his
Section~5.2, Tables~5.2, 5.3 and~5.5). BRAM in 36\,Kb blocks.}
\label{tab:kc_hwres}
\small
\begin{tabular}{@{}lr@{}}
\toprule
Quantity & Value\\
\midrule
Compression IP: LUT / FF / BRAM / DSP & \num{45892} / \num{9243} / 10 / 44\\
Whole system: LUT / FF / BRAM / DSP & \num{59353} / \num{32552} / 38 / 80\\
IP latency (HLS estimate) & 23 cycles, \SI{230}{\nano\second}\\
System latency (last DMA read to last DMA write) & \SI{79.4}{\micro\second}\\
Input throughput & ``\SI{11.92}{Gbit/s}''\\
CR: MATLAB / bit-accurate simulation / hardware & 3.25 / 3.24 / 3.17\\
\bottomrule
\end{tabular}
\end{table}

\section{What he left open}
\label{sec:kc_open}

He concluded that DRHE met every requirement of his use case. His future work
included floating-point arithmetic for the prediction, since his core far
exceeded its real-time requirement; a bit-accurate decompression model; tests
with a real sensor; safe error handling; and a guaranteed compression ratio.
This thesis takes up the first two: its DRHE predicts in floating point, and
Chapter~\ref{ch:hwinterpret} measures what that costs, and it implements and
verifies a hardware decompressor for every design. Its saturation of the one
residual the code cannot represent (Section~\ref{sec:robust}) is a form of the
safe error handling he asked for.
\label{ch:kiem:end}
